\section{AHSME 1979}
        \begin{enumerate}
        	\item (\textit{AHSME 1979}) 
                    \begin{figure}[H]
                        \centering
                        \includegraphics[width=0.5\textwidth]{Problem/AHSME_I_1979_1_Problem_Figure_1.png}
                    \end{figure}
                    


 If rectangle ABCD has area 72 square meters and E and G are the midpoints of sides AD and CD, respectively, then the area of rectangle DEFG in square meters is


 $\textbf{(A) }8\qquad \textbf{(B) }9\qquad \textbf{(C) }12\qquad \textbf{(D) }18\qquad \textbf{(E) }24$

 \textbf{Answer: }D



\textbf{Solution}

Solution by e\_power\_pi\_times\_i

 Since the dimensions of $DEFG$ are half of the dimensions of $ABCD$, the area of $DEFG$ is $\dfrac{1}{2}\cdot\dfrac{1}{2}$ of $ABCD$, so the area of $ABCD$ is $\dfrac{1}{4}\cdot72 = \boxed{\textbf{(D) } 18}$.

 


	\item (\textit{AHSME 1979}) For all non-zero real numbers $x$ and $y$ such that $x-y=xy, \frac{1}{x}-\frac{1}{y}$ equals


 $\textbf{(A) }\frac{1}{xy}\qquad \textbf{(B) }\frac{1}{x-y}\qquad \textbf{(C) }0\qquad \textbf{(D) }-1\qquad \textbf{(E) }y-x$

 \textbf{Answer: }D



\textbf{Solution 1}

Moving all variables to one side of the equation, we can use Simon's Favorite Factoring Trick to factor the equation into 
 \[(x+1)(y-1) = -1\] Plugging in $-1$ and $1$ as the $x$ and $y$ sides respectively, we get $x = -2$ and $y = 2$. Plugging this in to $\frac{1}{x}-\frac{1}{y}$ gives us $\boxed{-1}$ as our final answer.

 



\textbf{Solution 2}

Notice that we can do $\frac{x-y}{xy} = \frac{xy}{xy}$. We are left with $\frac{1}{y} - \frac{1}{x} = 1$. Multiply by $-1$ to achieve $\frac{1}{x} - \frac{1}{y} = \boxed{-1}$.

 


	\item (\textit{AHSME 1979}) 
                    \begin{figure}[H]
                        \centering
                        \includegraphics[width=0.5\textwidth]{Problem/AHSME_I_1979_3_Problem_Figure_1.png}
                    \end{figure}
                    


 In the adjoining figure, $ABCD$ is a square, $ABE$ is an equilateral triangle and point $E$ is outside square $ABCD$. 
What is the measure of $\measuredangle AED$ in degrees?


 $\textbf{(A) }10\qquad \textbf{(B) }12.5\qquad \textbf{(C) }15\qquad \textbf{(D) }20\qquad \textbf{(E) }25$

 \textbf{Answer: }C



\textbf{Solution}

Solution by e\_power\_pi\_times\_i

 Notice that $\measuredangle DAE = 90^\circ+60^\circ = 150^\circ$ and that $AD = AE$. Then triangle $ADE$ is isosceles, so $\measuredangle AED = \dfrac{180^\circ-150^\circ}{2} = \boxed{\textbf{(C) } 15}$.

 



\textbf{Solution 2 (Obnoxiously tedious)}

WLOG, let the side length of the square and the equilateral triangle be $1$. $\angle{DAE}=90^\circ+60^\circ=150^\circ$. Apply the law of cosines then the law of sines, we find that $\angle{AED}=15^\circ$. Select $\boxed{C}$.

 Please don't try this solution. This is just obnoxiously tedious and impractical.

 


	\item (\textit{AHSME 1979}) For all real numbers $x, x[x\{x(2-x)-4\}+10]+1=$
$\textbf{(A) }-x^4+2x^3+4x^2+10x+1\qquad \textbf{(B) }-x^4-2x^3+4x^2+10x+1\qquad\\ \textbf{(C) }-x^4-2x^3-4x^2+10x+1\qquad \textbf{(D) }-x^4-2x^3-4x^2-10x+1\qquad\\ \textbf{(E) }-x^4+2x^3-4x^2+10x+1$

 \textbf{Answer: }E


	\item (\textit{AHSME 1979}) Find the sum of the digits of the largest even three digit number (in base ten representation) 
which is not changed when its units and hundreds digits are interchanged.


 $\textbf{(A) }22\qquad \textbf{(B) }23\qquad \textbf{(C) }24\qquad \textbf{(D) }25\qquad \textbf{(E) }26$

 \textbf{Answer: }D



\textbf{Solution}

Solution by e\_power\_pi\_times\_i

 Since the number doesn't change when the units and hundreds digits are switched, the number must be of the form $aba$. We want to create the largest even $3$-digit number, so $a = 8$ and $b = 9$. The sum of the digits is $8+9+8 = \boxed{\textbf{(D) } 25}$.

 


	\item (\textit{AHSME 1979}) $\frac{3}{2}+\frac{5}{4}+\frac{9}{8}+\frac{17}{16}+\frac{33}{32}+\frac{65}{64}-7=$
$\textbf{(A) }-\frac{1}{64}\qquad \textbf{(B) }-\frac{1}{16}\qquad \textbf{(C) }0\qquad \textbf{(D) }\frac{1}{16}\qquad \textbf{(E) }\frac{1}{64}$

 \textbf{Answer: }A



\textbf{Solution}

Solution by e\_power\_pi\_times\_i

 Simplifying, we have $\frac{96}{64}+\frac{80}{64}+\frac{72}{64}+\frac{68}{64}+\frac{66}{64}+\frac{65}{64}-7$, which is $\frac{96+80+72+68+66+65}{64}-7 = \frac{447}{64}-7 = 6\frac{63}{64}-7 = \boxed{\textbf{(A) } -\frac{1}{64}}$.

 


	\item (\textit{AHSME 1979}) The square of an integer is called a perfect square. If $x$ is a perfect square, the next larger perfect square is


 $\textbf{(A) }x+1\qquad \textbf{(B) }x^2+1\qquad \textbf{(C) }x^2+2x+1\qquad \textbf{(D) }x^2+x\qquad \textbf{(E) }x+2\sqrt{x}+1$

 \textbf{Answer: }E



\textbf{Solution}

Solution by e\_power\_pi\_times\_i

 Since $x$ is a perfect square, denote $k$ such that $k^2 = x$. Then the next perfect square is $(k+1)^2 = k^2+2k+1$. Substituting back in the equation $k = \sqrt{x}$, the next square is $\boxed{\textbf{(E) }x+2\sqrt{x}+1}$.

 


	\item (\textit{AHSME 1979}) Find the area of the smallest region bounded by the graphs of $y=|x|$ and $x^2+y^2=4$.


 $\textbf{(A) }\frac{\pi}{4}\qquad \textbf{(B) }\frac{3\pi}{4}\qquad \textbf{(C) }\pi\qquad \textbf{(D) }\frac{3\pi}{2}\qquad \textbf{(E) }2\pi$

 \textbf{Answer: }C



\textbf{Solution}

Solution by e\_power\_pi\_times\_i

 The graph of $x^2+y^2 = 4$ is a circle with radius $2$ centered at the origin. The graph of $y=|x|$ is the combined graphs of $y=x$ and $y=-x$ with a nonnegative y. Because the arguments of $y=x$ and $y=-x$ are $135^\circ$ and $45^\circ$ respectively, the angle between the graphs of $y=x$ and $y=-x$ is $90^\circ$. Thus, the smallest region bounded by the graphs is $\frac{1}{4}\cdot2^2\cdot\pi = \boxed{\textbf{(C) } \pi}$.

 


	\item (\textit{AHSME 1979}) The product of $\sqrt[3]{4}$ and $\sqrt[4]{8}$ equals


 $\textbf{(A) }\sqrt[7]{12}\qquad \textbf{(B) }2\sqrt[7]{12}\qquad \textbf{(C) }\sqrt[7]{32}\qquad \textbf{(D) }\sqrt[12]{32}\qquad \textbf{(E) }2\sqrt[12]{32}$

 \textbf{Answer: }E



\textbf{Solution}

Solution by e\_power\_pi\_times\_i

 $\sqrt[3]{4}$ and $\sqrt[4]{8}$ can be expressed as $2^{\frac{2}{3}}$ and $2^{\frac{3}{4}}$, so their product is $2^{\frac{17}{12}} = \boxed{\textbf{(E) } 2\sqrt[12]{32}}$.

 


	\item (\textit{AHSME 1979}) If $P_1P_2P_3P_4P_5P_6$ is a regular hexagon whose apothem (distance from the center to midpoint of a side) is $2$, 
and $Q_i$ is the midpoint of side $P_iP_{i+1}$ for $i=1,2,3,4$, then the area of quadrilateral $Q_1Q_2Q_3Q_4$ is


 $\textbf{(A) }6\qquad \textbf{(B) }2\sqrt{6}\qquad \textbf{(C) }\frac{8\sqrt{3}}{3}\qquad \textbf{(D) }3\sqrt{3}\qquad \textbf{(E) }4\sqrt{3}$

 \textbf{Answer: }D



\textbf{Solution}

Solution by e\_power\_pi\_times\_i

 Notice that quadrilateral $Q_1Q_2Q_3Q_4$ consists of $3$ equilateral triangles with side length $2$. Thus the area of the quadrilateral is $3\cdot(\frac{2^2\cdot\sqrt{3}}{4}) = 3\cdot\sqrt{3} = \boxed{\textbf{(D) } 3\sqrt{3}}$.

 


	\item (\textit{AHSME 1979}) Find a positive integral solution to the equation
$\frac{1+3+5+\dots+(2n-1)}{2+4+6+\dots+2n}=\frac{115}{116}$
$\textbf{(A) }110\qquad \textbf{(B) }115\qquad \textbf{(C) }116\qquad \textbf{(D) }231\qquad\\ \textbf{(E) }\text{The equation has no positive integral solutions.}$

 \textbf{Answer: }B



\textbf{Solution}

Solution by e\_power\_pi\_times\_i

 Notice that the numerator and denominator are the sum of the first $n$ odd and even numbers, respectively. Then the numerator is $n^2$, and the denominator is $n(n+1)$.  Then $\frac{n}{n+1} = \frac{115}{116}$, so $n = \boxed{\textbf{(B) } 115}$.

 


	\item (\textit{AHSME 1979}) 
                    \begin{figure}[H]
                        \centering
                        \includegraphics[width=0.5\textwidth]{Problem/AHSME_I_1979_12_Problem_Figure_1.png}
                    \end{figure}
                    


 In the adjoining figure, $CD$ is the diameter of a semi-circle with center $O$. Point $A$ lies on the extension of $DC$ past $C$; 
point $E$ lies on the semi-circle, and $B$ is the point of intersection (distinct from $E$ ) of line segment $AE$ with the semi-circle. 
If length $AB$ equals length $OD$, and the measure of $\measuredangle EOD$ is $45^\circ$, then the
measure of $\measuredangle BAO$ is


 $\textbf{(A) }10^\circ\qquad \textbf{(B) }15^\circ\qquad \textbf{(C) }20^\circ\qquad \textbf{(D) }25^\circ\qquad \textbf{(E) }30^\circ$

 \textbf{Answer: }B



\textbf{Solution}

Solution by e\_power\_pi\_times\_i

 Because $AB = OD$, triangles $ABO$ and $BOE$ are isosceles. Denote $\measuredangle BAO = \measuredangle AOB = \theta$. Then $\measuredangle ABO = 180^\circ-2\theta$, and $\measuredangle EBO = \measuredangle OEB = 2\theta$, so $\measuredangle BOE = 180^\circ-4\theta$. Notice that $\measuredangle AOB + \measuredangle BOE + 45^\circ = 180^\circ$. Therefore $\theta+180-4\theta = 135^\circ$, and $\theta = \boxed{\textbf{(B) } 15^\circ}$.

 



\textbf{Solution 2}

Draw $BO$. Let $y = \angle BAO$. Since $AB = OD = BO$, triangle $ABO$ is isosceles, so $\angle BOA = \angle BAO = y$. Angle $\angle EBO$ is exterior to triangle $ABO$, so $\angle EBO = \angle BAO + \angle BOA = y + y = 2y$.

 Triangle $BEO$ is isosceles, so $\angle BEO = \angle EBO = 2y$. Then $\angle EOD$ is external to triangle $AEO$, so $\angle EOD = \angle EAO + \angle AEO = y + 2y = 3y$. But $\angle EOD = 45^\circ$, so $\angle BAO = y = 45^\circ/3 = \boxed{15^\circ}$. That means the answer is $\boxed{\textbf{(B) } 15^\circ}$.

 


	\item (\textit{AHSME 1979}) The inequality $y-x<\sqrt{x^2}$ is satisfied if and only if


 $\textbf{(A) }y<0\text{ or }y<2x\text{ (or both inequalities hold)}\qquad \textbf{(B) }y>0\text{ or }y<2x\text{ (or both inequalities hold)}\qquad \textbf{(C) }y^2<2xy\qquad \textbf{(D) }y<0\qquad \textbf{(E) }x>0\text{ and }y<2x$

 \textbf{Answer: }A



\textbf{Solution}

Solution by e\_power\_pi\_times\_i

 $\sqrt{x^2} = \pm x$, so the inequality is just $y-x<\pm x$. Therefore we get the two inequalities $y<0$ and $y<2x$. Checking the answer choices, we find that $\boxed{\textbf{(A) } y<0\text{ or }y<2x\text{ (or both inequalities hold)}}$ is the answer.

 


	\item (\textit{AHSME 1979}) In a certain sequence of numbers, the first number is $1$, and, for all $n\ge 2$, the product of the first $n$ numbers in the sequence is $n^2$. 
The sum of the third and the fifth numbers in the sequence is


 $\textbf{(A) }\frac{25}{9}\qquad \textbf{(B) }\frac{31}{15}\qquad \textbf{(C) }\frac{61}{16}\qquad \textbf{(D) }\frac{576}{225}\qquad \textbf{(E) }34$

 \textbf{Answer: }C



\textbf{Solution}

Solution by e\_power\_pi\_times\_i

 Since the product of the first $n$ numbers in the sequence is $n^2$, the product of the first $n+1$ numbers in the sequence is $(n+1)^2$. Therefore the $n+1$th number in the sequence is $\frac{(n+1)^2}{n^2}$. Therefore the third and the fifth numbers are $\frac{9}{4}$ and $\frac{25}{16}$ respectively. The sum of those numbers is $\frac{36}{16}+\frac{25}{16} = \boxed{\textbf{(C) } \frac{61}{16}}$.

 


	\item (\textit{AHSME 1979}) Two identical jars are filled with alcohol solutions, the ratio of the volume of alcohol to the volume of water being $p : 1$ 
in one jar and $q : 1$ in the other jar. If the entire contents of the two jars are mixed together, 
the ratio of the volume of alcohol to the volume of water in the mixture is


 $\textbf{(A) }\frac{p+q}{2}\qquad \textbf{(B) }\frac{p^2+q^2}{p+q}\qquad \textbf{(C) }\frac{2pq}{p+q}\qquad \textbf{(D) }\frac{2(p^2+pq+q^2)}{3(p+q)}\qquad \textbf{(E) }\frac{p+q+2pq}{p+q+2}$

 \textbf{Answer: }E



\textbf{Solution}

Solution by e\_power\_pi\_times\_i

 The amount of alcohol in the jars are $\frac{p}{p+1}$ and $\frac{q}{q+1}$, and the amount of water in the jars are $\frac{1}{p+1}$ and $\frac{1}{q+1}$. Then the total amount of alcohol is $\frac{p}{p+1} + \frac{q}{q+1} = \frac{p+q+2pq}{(p+1)(q+1)}$, and the total amount of water is $\frac{1}{p+1} and \frac{1}{q+1} = \frac{p+q+2}{(p+1)(q+1)}$. The ratio of the volume of alcohol to the volume of water in the mixture is $\frac{\frac{p+q+2pq}{(p+1)(q+1)}}{\frac{p+q+2}{(p+1)(q+1)}} = \boxed{\textbf{(E) } \frac{p+q+2pq}{p+q+2}}$.

 


	\item (\textit{AHSME 1979}) A circle with area $A_1$ is contained in the interior of a larger circle with area $A_1+A_2$. If the radius of the larger circle is $3$, 
and if $A_1 , A_2, A_1 + A_2$ is an arithmetic progression, then the radius of the smaller circle is


 $\textbf{(A) }\frac{\sqrt{3}}{2}\qquad \textbf{(B) }1\qquad \textbf{(C) }\frac{2}{\sqrt{3}}\qquad \textbf{(D) }\frac{3}{2}\qquad \textbf{(E) }\sqrt{3}$

 \textbf{Answer: }E



\textbf{Solution}

Solution by e\_power\_pi\_times\_i

 The area of the larger circle is $A_1 + A_2 = 9\pi$. Then $A_1 , 9\pi-A_1 , 9\pi$ are in an arithmetic progression. Thus $9\pi-(9\pi-A_1) = 9\pi-A_1-A_1$. This simplifies to $3A_1 = 9\pi$, or $A_1 = 3\pi$. The radius of the smaller circle is $\boxed{\textbf{(E) } \sqrt{3}}$.

 


	\item (\textit{AHSME 1979}) 
                    \begin{figure}[H]
                        \centering
                        \includegraphics[width=0.5\textwidth]{Problem/AHSME_I_1979_17_Problem_Figure_1.png}
                    \end{figure}
                    


 Points $A , B, C$, and $D$ are distinct and lie, in the given order, on a straight line. 
Line segments $AB, AC$, and $AD$ have lengths $x, y$, and $z$, respectively. 
If line segments $AB$ and $CD$ may be rotated about points $B$ and $C$, respectively, 
so that points $A$ and $D$ coincide, to form a triangle with positive area, 
then which of the following three inequalities must be satisfied?


 $\textbf{I. }x<\frac{z}{2}\qquad\\ \textbf{II. }y<x+\frac{z}{2}\qquad\\ \textbf{III. }y<\frac{z}{2}\qquad$
$\textbf{(A) }\textbf{I. }\text{only}\qquad \textbf{(B) }\textbf{II. }\text{only}\qquad \textbf{(C) }\textbf{I. }\text{and }\textbf{II. }\text{only}\qquad \textbf{(D) }\textbf{II. }\text{and }\textbf{III. }\text{only}\qquad \textbf{(E) }\textbf{I. },\textbf{II. },\text{and }\textbf{III. }$

 \textbf{Answer: }C



\textbf{Solution}

Solution by e\_power\_pi\_times\_i

 We know that this triangle has lengths of $x$, $y-x$, and $z-y$. Using the Triangle Inequality, we get $3$ inequalities: $2y>z, z>2x, 2x+z>0$. Therefore, we know that $\textbf{I}$ is true and $\textbf{III}$ is false. In $\textbf{II}$, we have to prove $2y<2x+z$. We know that $2y>z$, so we have to prove $z<2y<2x+z$. $2x<z$, so we have to prove that $z<2z$, which is true for all positive $z$. Therefore the answer is $\boxed{\textbf{(C) } \textbf{I. }\text{and }\textbf{II. }\text{only}}$.

 


	\item (\textit{AHSME 1979}) To the nearest thousandth, $\log_{10}2$ is $.301$ and $\log_{10}3$ is $.477$. 
Which of the following is the best approximation of $\log_5 10$?


 $\textbf{(A) }\frac{8}{7}\qquad \textbf{(B) }\frac{9}{7}\qquad \textbf{(C) }\frac{10}{7}\qquad \textbf{(D) }\frac{11}{7}\qquad \textbf{(E) }\frac{12}{7}$

 \textbf{Answer: }C



\textbf{Solution}

Solution by e\_power\_pi\_times\_i

 Notice that $\log_5 10 = \log_5 2 + 1 = \frac{1}{\log_2 5} + 1$. So we are trying to find $\log_2 5$. Denote $\log_{10}2$ as $x$. Then $\frac{1}{x} = \log_2 10 = \log_2 5 + 1$. Therefore $\log_2 5 = \frac{1}{x}-1$, and plugging this in gives $\log_5 10 = \frac{1}{\frac{1}{x}-1}+1 = \frac{x}{1-x}+1 = \frac{1}{1-x}$. Since $x$ is around $\frac{3}{10}$, we substitute and get $\boxed{\textbf{(C) }\frac{10}{7}}$.

 


	\item (\textit{AHSME 1979}) Find the sum of the squares of all real numbers satisfying the equation $x^{256}-256^{32}=0$.


 $\textbf{(A) }8\qquad \textbf{(B) }128\qquad \textbf{(C) }512\qquad \textbf{(D) }65,536\qquad \textbf{(E) }2(256^{32})$

 \textbf{Answer: }A



\textbf{Solution}

Solution by e\_power\_pi\_times\_i

 Notice that the solutions to the equation $x^{256}-1=0$ are the $256$ roots of unity. Then the solutions to the equation $x^{256}-256^{32}=0$ are the $256$ roots of unity dilated by $\sqrt[256]{256^{32}} = \sqrt[256]{2^{256}} = 2$. However, the only real solutions to the equation are the first root of unity and the root of unity opposite of it, as both are on the real axis in the complex plane. These two roots of unity are $\pm1$, and dilating by $2$ gives $\pm2$. The sum of the squares is $(2)^2+(-2)^2 = \boxed{\textbf{(A) } 8}$




\textbf{Solution 2}

Notice that we can change the equation to be $x^{256}=256^{32}$.

 The RHS can be simplified to $x^{256}=2^{256}$

 Hence, the only real solutions are $x=-2,2$

 Hence, $(-2)^2+(2)^2=8=\boxed{A}$.

 


	\item (\textit{AHSME 1979}) If $a=\frac{1}{2}$ and $(a+1)(b+1)=2$ then the radian measure of $\arctan a + \arctan b$ equals


 $\textbf{(A) }\frac{\pi}{2}\qquad \textbf{(B) }\frac{\pi}{3}\qquad \textbf{(C) }\frac{\pi}{4}\qquad \textbf{(D) }\frac{\pi}{5}\qquad \textbf{(E) }\frac{\pi}{6}$

 \textbf{Answer: }C



\textbf{Solution}

Solution by e\_power\_pi\_times\_i

 Since $a=\frac{1}{2}$, $b=\frac{1}{3}$. Now we evaluate $\arctan a$ and $\arctan b$. Denote $x$ and $\theta$ such that $\arctan x = \theta$. Then $\tan(\arctan(x)) = \tan(\theta)$, and simplifying gives $x = \tan(\theta)$. So $a = \tan(\theta_a) = \frac{1}{2}$ and $b = \tan(\theta_b) = \frac{1}{3}$. The question asks for $\theta_a + \theta_b$, so we try to find $\tan(\theta_a + \theta_b)$ in terms of $\tan(\theta_a)$ and $\tan(\theta_b)$. Using the angle addition formula for $\tan(\alpha+\beta)$, we get that $\tan(\theta_a + \theta_b) = \frac{\tan(\theta_a)+\tan(\theta_b)}{1-\tan(\theta_a)\tan(\theta_b)}$. Plugging $\tan(\theta_a) = \frac{1}{2}$ and $\tan(\theta_b) = \frac{1}{3}$ in, we have $\tan(\theta_a + \theta_b) = \frac{\frac{1}{2}+\frac{1}{3}}{1-(\frac{1}{2})(\frac{1}{3})}$. Simplifying, $\tan(\theta_a + \theta_b) = 1$, so $\theta_a + \theta_b$ in radians is $\boxed{\textbf{(C) } \frac{\pi}{4}}$.

 



\textbf{Solution 1.1 (Guiding through the thought process)}

Thinking through the problem, we can see $\arctan a$ is the angle whose tangent is $0.5$. $\arctan b$ is the angle whose tangent is $\frac{1}{3}$. Call the former angle $\theta_1$, the latter $\theta_2$. So we are trying to find $\theta_1+\theta_2$. So given two tangent measures, it is natural for us to think about the sum of tangent measures (what else can we try? Remember: we are not allowed to use calculators). Plug in as above and continue on.

 


	\item (\textit{AHSME 1979}) The length of the hypotenuse of a right triangle is $h$ , and the radius of the inscribed circle is $r$. 
The ratio of the area of the circle to the area of the triangle is


 $\textbf{(A) }\frac{\pi r}{h+2r}\qquad \textbf{(B) }\frac{\pi r}{h+r}\qquad \textbf{(C) }\frac{\pi}{2h+r}\qquad \textbf{(D) }\frac{\pi r^2}{r^2+h^2}\qquad \textbf{(E) }\text{none of these}$

 \textbf{Answer: }B



\textbf{Solution}

Solution by e\_power\_pi\_times\_i

 Since $rs = A$, where $r$ is the inradius, $s$ is the semiperimeter, and $A$ is the area, we have that the ratio of the area of the circle to the area of the triangle is $\frac{\pi r^2}{rs} = \frac{\pi r}{s}$. Now we try to express $s$ as $h$ and $r$. Denote the points where the incircle meets the triangle as $X,Y,Z$, where $O$ is the incenter, and denote $AX = AY = z, BX = BZ = y, CY = CZ = x$. Since $XOZB$ is a square (tangents are perpendicular to radius), $r = BX = BZ = y$. The perimeter can be expressed as $2(x+y+z)$, so the semiperimeter is $x+y+z$. The hypotenuse is $AY+CY = z+x$. Thus we have $s = x+y+z = (z+x)+y = h+r$. The answer is $\boxed{\textbf{(B) } \frac{\pi r}{h+r}}$.

 


	\item (\textit{AHSME 1979}) Find the number of pairs $(m, n)$ of integers which satisfy the equation $m^3 + 6m^2 + 5m = 27n^3 + 9n^2 + 9n + 1$.


 $\textbf{(A) }0\qquad \textbf{(B) }1\qquad \textbf{(C) }3\qquad \textbf{(D) }9\qquad \textbf{(E) }\infty$

 \textbf{Answer: }A



\textbf{Solution}

The equation is equivalent to $m(m+1)(m+5) = 27n^3 + 27n^2 + 9n + 1$. Taking mod 3, we get$m(m+1)(m+2)=1 (\bmod 3)$. However, $m(m+1)(m+2)$ is always divisible by $3$ for any integer $m$. Thus, the answer is $\boxed{\textbf{(A)} 0}$Solution by mickyboy789

 


	\item (\textit{AHSME 1979}) The edges of a regular tetrahedron with vertices $A ,~ B,~ C$, and $D$ each have length one. 
Find the least possible distance between a pair of points $P$ and $Q$, where $P$ is on edge $AB$ and $Q$ is on edge $CD$.


 


                    \begin{figure}[H]
                        \centering
                        \includegraphics[width=0.5\textwidth]{Problem/AHSME_I_1979_23_Problem_Figure_1.png}
                    \end{figure}
                    
$\textbf{(A) }\frac{1}{2}\qquad \textbf{(B) }\frac{3}{4}\qquad \textbf{(C) }\frac{\sqrt{2}}{2}\qquad \textbf{(D) }\frac{\sqrt{3}}{2}\qquad \textbf{(E) }\frac{\sqrt{3}}{3}$

 \textbf{Answer: }C



\textbf{Solution 1}

Note that the distance $PQ$ will be minimized when $P$ is the midpoint of $AB$ and $Q$ is the midpoint of $CD$.

 To find this distance, consider triangle $\triangle PCQ$. $Q$ is the midpoint of $CD$, so $CQ=\frac{1}{2}$. Additionally, since $CP$ is the altitude of equilateral $\triangle ABC$, $CP=\frac{\sqrt{3}}{2}$.

 
                    \begin{figure}[H]
                        \centering
                        \includegraphics[width=0.5\textwidth]{Solution/AHSME_I_1979_23_Solution_Figure_1.png}
                    \end{figure}
                    

 Next, we need to find $\cos(\angle PCQ)$ in order to find $PQ$ by the Law of Cosines. To do so, drop down $D$ onto $\triangle ABC$ to get the point $D^\prime$.

 $\angle PCD$ is congruent to $\angle D^\prime CD$, since $P$, $D^\prime$, and $C$ are collinear. Therefore, we can just find $\cos(\angle D^\prime CD)$.

 Note that $\triangle CD^\prime D$ is a right triangle with $\angle CD^\prime D$ as a right angle. 

 

                    \begin{figure}[H]
                        \centering
                        \includegraphics[width=0.5\textwidth]{Solution/AHSME_I_1979_23_Solution_Figure_2.png}
                    \end{figure}
                    

 
As given by the problem, $CD=1$.

 Note that $D^\prime$ is the centroid of equilateral $\triangle ABC$. Additionally, since $\triangle ABC$ is equilateral, $D^\prime$ is also the orthocenter. Due to this, the distance from $C$ to $D^\prime$ is $\frac{2}{3}$ of the altitude of $\triangle ABC$. Therefore, $CD^\prime=\frac{\sqrt{3}}{3}$. 

 Since $\cos(\angle D^\prime CD)=\cos(\angle PCQ)=\frac{CD^\prime}{CD}$, $\cos(\angle PCQ)=\frac{\frac{\sqrt{3}}{3}}{1}=\frac{\sqrt{3}}{3}$

 \[PQ^2=CP^2+CQ^2-2(CP)(CQ)\cos(\angle PCQ)\]
 \[PQ^2=\frac{3}{4}+\frac{1}{4}-2\left(\frac{\sqrt{3}}{4}\right)\left(\frac{1}{2}\right)\left(\frac{\sqrt{3}}{3}\right)\]Simplifying, $PQ^2=\frac{1}{2}$.Therefore, $PQ=\frac{\sqrt{2}}{2}\Rightarrow$ $\boxed{\textbf{C}}$

 Solution by treetor10145

 



\textbf{Solution 2 (less overkill)}

Notice, like above said, that $P$ is the midpoint of $AB$ and $Q$ is the midpoint of $CD$. 

 
To find the length of $PQ$, first draw in lines $CP$ and $DP$. Notice that $DP$ is an altitude of $\triangle ADP$. We find that $\angle{DAP} = 60 ^{\circ}$ (since $\triangle ABD$ is equilateral), and $AD=\frac{1}{2}$. Use the properties of 30-60-90 triangles to get $DP=\frac{\sqrt{3}}{2}$. Since $CP$ is an altitude of a congruent equilateral triangle, $CP=DP=\frac{\sqrt{3}}{2}$.

 


 Notice that $\triangle CDP$ is isosceles with $CP=DP$. Also, since $Q$ is the midpoint of base $CD$, we can conclude that $PQ$ is an altitude. We can use Pythagorean theorem to get the following (taking into consideration $DQ=\frac{1}{2}$):

 


 
 \[DQ^2+PQ^2=PD^2\]

 \[\left(\frac{1}{2}\right)^2+PQ^2 = \left(\frac{\sqrt{3}}{2}\right)^2\]

 \[PQ^2=\frac{3}{4}-\frac{1}{4}=\frac{1}{2}\]

 \[PQ=\frac{\sqrt{2}}{2}\Rightarrow \boxed{\textbf{C}}\]

 
-WannabeCharmander

 


	\item (\textit{AHSME 1979}) Sides $AB,~ BC$, and $CD$ of (simple*) quadrilateral $ABCD$ have lengths $4,~ 5$, and $20$, respectively. 
If vertex angles $B$ and $C$ are obtuse and $\sin C = - \cos B =\frac{3}{5}$, then side $AD$ has length


 
            \begin{itemize}
        \item A polygon is called “simple” if it is not self intersecting.
            \end{itemize}
            $\textbf{(A) }24\qquad \textbf{(B) }24.5\qquad \textbf{(C) }24.6\qquad \textbf{(D) }24.8\qquad \textbf{(E) }25$

 \textbf{Answer: }E



\textbf{Solution}

We know that $\sin(C)=-\cos(B)=\frac{3}{5}$. Since $B$ and $C$ are obtuse, we have $\sin(180-C)=\cos(180-B)=\frac{3}{5}$. It is known that $\sin(x)=\cos(90-x)$, so $180-C=90-(180-C)=180-B$. We simplify this as follows:

 
 \[-90+C=180-B\]

 \[B+C=270^{\circ}\]

 Since $B+C=270^{\circ}$, we know that $A+D=360-(B+C)=90^{\circ}$. Now extend $AB$ and $CD$ as follows:

 
                    \begin{figure}[H]
                        \centering
                        \includegraphics[width=0.5\textwidth]{Solution/AHSME_I_1979_24_Solution_Figure_1.png}
                    \end{figure}
                    

 Let $AB$ and $CD$ intersect at $E$. We know that $\angle AED=90^{\circ}$ because $\angle E = 180 - (A+D)=180-90 = 90^{\circ}$.

 Since $\sin BCD = \frac{3}{5}$, we get $\sin ECB=\sin(180-BCD)=\sin BCD = \frac{3}{5}$. Thus, $EB=3$ and $EC=4$ from simple sin application.

 $AD$ is the hypotenuse of right $\triangle AED$, with leg lengths $AB+BE=7$ and $EC+CD=24$. Thus, $AD=\boxed{\textbf{(E)}25}$

 -WannabeCharmander

 


	\item (\textit{AHSME 1979}) If $q_1 ( x )$ and $r_ 1$ are the quotient and remainder, respectively, when the polynomial $x^ 8$ is divided by 
$x + \frac{1}{2}$ , and if $q_ 2 ( x )$ and $r_2$ are the quotient and remainder, respectively, 
when $q_ 1 ( x )$ is divided by $x + \frac{1}{2}$, then find $r_2$.


 $\textbf{(A) }\frac{1}{256}\qquad \textbf{(B) }-\frac{1}{16}\qquad \textbf{(C) }1\qquad \textbf{(D) }-16\qquad \textbf{(E) }256$

 \textbf{Answer: }B



\textbf{Solution}

Solution by e\_power\_pi\_times\_i

 First, we divide $x^8$ by $x+\frac{1}{2}$ using synthetic division or some other method. The quotient is $x^7-\frac{1}{2}x^6+\frac{1}{4}x^5-\frac{1}{8}x^4+\frac{1}{16}x^3-\frac{1}{32}x^2+\frac{1}{64}x-\frac{1}{128}$, and the remainder is $\frac{1}{256}$. Then we plug the solution to $x+\frac{1}{2} = 0$ into the quotient to find the remainder. Notice that every term in the quotient, when $x=-\frac{1}{2}$, evaluates to $-\frac{1}{128}$. Thus $r_2 =-\frac{8}{128} = \boxed{\textbf{(B) } -\frac{1}{16}}$.

 



\textbf{Solution 2}

Using the remainder theorem, we see that the remainder upon dividing $x^8$ by $x+\frac{1}{2}$ is equal to $\left(-\frac{1}{2}\right)^8 = \frac{1}{256}.$ Thus, we have that
 \[x^8 = \left(x+\frac{1}{2}\right)Q(x) + \frac{1}{256}.\]Isolating $Q(x),$ we see
 \[\frac{x^8-\frac{1}{256}}{x+\frac{1}{2}} = Q(x).\]Using the difference of squares factorization repeatedly, we get
 \[\frac{\left(x-\frac{1}{2}\right)\left(x+\frac{1}{2}\right)\left(x^2+\frac{1}{2^2}\right)\left(x^4+\frac{1}{2^4}\right)}{x+\frac{1}{2}}=\left(x-\frac{1}{2}\right)\left(x^2+\frac{1}{2^2}\right)\left(x^4+\frac{1}{2^4}\right)=Q(x).\]Finally, plugging in $x=-\frac{1}{2}$ again, the final answer is
 \[(-1)\left(\frac{1}{2}\right)\left(\frac{1}{8}\right)\]
 \[=\boxed{-\frac{1}{16}}.\]



	\item (\textit{AHSME 1979}) The function $f$ satisfies the functional equation $f(x) +f(y) = f(x + y ) - xy - 1$ for every pair $x,~ y$
of real numbers. If $f( 1) = 1$, then the number of integers $n \neq 1$ for which $f ( n ) = n$ is


 $\textbf{(A) }0\qquad \textbf{(B) }1\qquad \textbf{(C) }2\qquad \textbf{(D) }3\qquad \textbf{(E) }\infty$

 \textbf{Answer: }B



\textbf{Solution}

We are given that $f(x) +f(y) = f(x + y ) - xy - 1$ and $f( 1) = 1$, so we can let $y = 1$.  Thus we have:

 
 \[f(x) + 1 = f(x + 1) - x - 1\]

 Rearranging gives a recursive formula for $f$:

 
 \[f(x + 1) = f(x) + x + 2\]

 We notice that this is the recursive form for a quadratic, so f(x) must be of the form $ax^{2} + bx + c$.  To solve for $a, b,$ and $c$, we can first work backwards to solve for the values of f(0) and f(-1):

 
 \[f(1) = f(0) + 0 + 2 \Rightarrow 1 = f(0) + 2 \Rightarrow f(0) = -1\]
 \[f(0) = f(-1)  -1 + 2 \Rightarrow -1 = f(-1) + 1 \Rightarrow f(-1) = -2\]

 Since $f(0) = -1$:
 \[f(0) = a(0)^{2} + b(0) + c = -1 \Rightarrow c = -1\]  Since $f(1) = 1$: 
 \[f(1) = a(1)^{2} + b(1)+ c = a + b - 1 = 1 \Rightarrow a + b = 2\] Similarly, since $f(-1) = -2$:
 \[f(-1) = a(-1)^{2} + b(-1)+ c = a - b - 1 = -2 \Rightarrow a - b = -1\]

 Thus we have the system of equations:

 
 \[a + b = 2\]
 \[a - b = -1\]

 Which can be solved to yield $a = \frac{1}{2}$, $b = \frac{3}{2}$.  Therefore, $f(x) = \frac{1}{2} x^{2} + \frac{3}{2} x - 1$.

 Since we are searching for values for which $f(x) = x$, we have the equation $\frac{1}{2} x^{2} + \frac{3}{2} x - 1 = x$.  Subtracting $x$ yields $\frac{1}{2} x^{2} + \frac{1}{2} x - 1 = 0$, which we can simplify by dividing both sides by $\frac{1}{2}$: $x^{2} + x - 2 = 0$.  This factors into $(x + 2) (x - 1) = 0$, so therefore there are two solutions to $f(x) = x$: $-2$ and $1$.  Since the problem asks only for solutions that do not equal $1$, the answer is $\fbox{(B) 1}$.

 


	\item (\textit{AHSME 1979}) An ordered pair $( b , c )$ of integers, each of which has absolute value less than or equal to five, is chosen at random, with each
such ordered pair having an equal likelihood of being chosen. What is the probability that the equation $x^ 2 + bx + c = 0$ will
not have distinct positive real roots?


 $\textbf{(A) }\frac{106}{121}\qquad \textbf{(B) }\frac{108}{121}\qquad \textbf{(C) }\frac{110}{121}\qquad \textbf{(D) }\frac{112}{121}\qquad \textbf{(E) }\text{none of these}$

 \textbf{Answer: }E



\textbf{Solution}

$\boxed{E}$ 

 There are $10$ cases where the roots are real, positive and distinct.The roots to the quadratic equation are $\frac{-b\pm\sqrt{b^2-4c}}{2}$. 

 In order for $-b-\sqrt{b^2-4c}>0$ we certainly need $b<0$.Additionally, for the roots to be both real and distinct we need $b^2-4c>0$. And lastly, for both roots to be positive weneed $-b>\sqrt{b^2-4c}$. Combining these inequalities we determine that  $\frac{b^2}{4}>c>0$ and $b<0$. 

 This leaves just $10$ cases $\{(-5,1),(-5,2)(-5,3),(-5,4),(-5,5),(-4,1),(-4,2),(-4,3),(-3,1),(-3,2)\}$ which can easily be checked by hand.

 
The answer is $\frac{111}{121}$



	\item (\textit{AHSME 1979}) 
                    \begin{figure}[H]
                        \centering
                        \includegraphics[width=0.5\textwidth]{Problem/AHSME_I_1979_28_Problem_Figure_1.png}
                    \end{figure}
                    


 Circles with centers $A ,~ B$, and $C$ each have radius $r$, where $1 < r < 2$. 
The distance between each pair of centers is $2$. If $B'$ is the point of intersection of circle $A$ and circle $C$ 
which is outside circle $B$, and if $C'$ is the point of intersection of circle $A$ and circle $B$ which is outside circle $C$, 
then length $B'C'$ equals


 $\textbf{(A) }3r-2\qquad \textbf{(B) }r^2\qquad \textbf{(C) }r+\sqrt{3(r-1)}\qquad\\ \textbf{(D) }1+\sqrt{3(r^2-1)}\qquad \textbf{(E) }\text{none of these}$

 \textbf{Answer: }D



\textbf{Solution 1 (Coordinate Geometry)}

The circles can be described in the cartesian plane as being centered at $(-1,0),(1,0)$ and $(0,\sqrt{3})$ with radius $r$ by the equations 

 $x^2+(y-\sqrt{3})^2=r^2$
$(x+1)^2+y^2=r^2$
$(x-1)^2+y^2=r^2$.

 Solving the first 2 equations gives $x=1-\sqrt{3}\cdot y$ which when substituted back in gives $y=\frac{\sqrt{3}\pm \sqrt{r^2-1}}{2}$.

 The larger root $y=\frac{\sqrt{3}+\sqrt{r^2-1}}{2}$ is the point B' described in the question. This root corresponds to $x=-\frac{1+\sqrt{3(r^2-1)}}{2}$. 

 By symmetry across the y-axis the length of the line segment $B'C'$ is $1+\sqrt{3(r^2-1)}$ which is $\boxed{D}$.

 



\textbf{Solution 2 (Synthetic)}

Suppose $B’B$ and $AC$ intersect at $P$. By the Pythagorean Theorem, $B’P = \sqrt{r^2 - 1}$ and by a $30-60-90$ triangle, $PB = \sqrt{3}$. Using Ptolemy’s Theorem on isosceles trapezoid $BCB’C’$, we get that 
 \[2(B’C’) + r^2 = (\sqrt{3} + \sqrt{r^2 - 1})^2.\] After a little algebra, we get that $B’C’ = \boxed{1 + \sqrt{3(r^2 - 1)}}$ as desired.\href{https://artofproblemsolving.com/wiki/index.php?title=User:Solasky}{Solasky} (\href{https://artofproblemsolving.com/wiki/index.php?title=User\_talk:Solasky&action=edit&redlink=1}{talk}) 12:29, 27 May 2023 (EDT)

 


	\item (\textit{AHSME 1979}) For each positive number $x$, let 
$f(x)=\frac{\left(x+\frac{1}{x}\right)^6-\left(x^6+\frac{1}{x^6}\right)-2} {\left(x+\frac{1}{x}\right)^3+\left(x^3+\frac{1}{x^3}\right)}$. 
The minimum value of $f(x)$ is


 $\textbf{(A) }1\qquad \textbf{(B) }2\qquad \textbf{(C) }3\qquad \textbf{(D) }4\qquad \textbf{(E) }6$

 \textbf{Answer: }E



\textbf{Solution}

Let $a = \left( x + \frac{1}{x} \right)^3$ and $b = x^3 + \frac{1}{x^3}$. Then
 \begin{align*} f(x) &= \frac{\left( x + \frac{1}{x} \right)^6 - (x^6 + \frac{1}{x^6}) - 2}{\left( x + \frac{1}{x} \right)^3 + (x^3 + \frac{1}{x^3})} \\ &= \frac{\left( x + \frac{1}{x} \right)^6 - (x^6 + 2 + \frac{1}{x^6})}{\left( x + \frac{1}{x} \right)^3 + (x^3 + \frac{1}{x^3})} \\ &= \frac{\left( x + \frac{1}{x} \right)^6 - (x^3 + \frac{1}{x^3})^2}{\left( x + \frac{1}{x} \right)^3 + (x^3 + \frac{1}{x^3})} \\ &= \frac{a^2 - b^2}{a + b}. \end{align*}


 By difference of squares,
 \begin{align*} f(x) &= \frac{(a - b)(a + b)}{a + b} \\ &= a - b \\ &= \left( x + \frac{1}{x} \right)^3 - \left( x^3 + \frac{1}{x^3} \right) \\ &= \left( x^3 + 3x + \frac{3}{x} + \frac{1}{x^3} \right) - \left( x^3 + \frac{1}{x^3} \right) \\ &= 3x + \frac{3}{x} \\ &= 3 \left( x + \frac{1}{x} \right). \end{align*}


 By the AM-GM inequality,
 \[x + \frac{1}{x} \ge 2,\]so $f(x) \ge 6$. Furthermore, when $x = 1$, $f(1) = 6$, so the minimum value of $f(x)$ is $\boxed{6}$.

 


	\item (\textit{AHSME 1979}) 
                    \begin{figure}[H]
                        \centering
                        \includegraphics[width=0.5\textwidth]{Problem/AHSME_I_1979_30_Problem_Figure_1.png}
                    \end{figure}
                    


 In $\triangle ABC$, $E$ is the midpoint of side $BC$ and $D$ is on side $AC$. 
If the length of $AC$ is $1$ and $\measuredangle BAC = 60^\circ, \measuredangle ABC = 100^\circ, \measuredangle ACB = 20^\circ$ and 
$\measuredangle DEC = 80^\circ$, then the area of $\triangle ABC$ plus twice the area of $\triangle CDE$ equals


 $\textbf{(A) }\frac{1}{4}\cos 10^\circ\qquad \textbf{(B) }\frac{\sqrt{3}}{8}\qquad \textbf{(C) }\frac{1}{4}\cos 40^\circ\qquad \textbf{(D) }\frac{1}{4}\cos 50^\circ\qquad \textbf{(E) }\frac{1}{8}$

 \textbf{Answer: }B



\textbf{Solution}

Let $F$ be the point on the extension of side $AB$ past $B$ for which $AF=1$. Since $AF=AC$ and $\measuredangle FAC = 60^\circ$,$\triangle ACF$ is equilateral. Let $G$ be the point on line segment $BF$ for which $\measuredangle BCG=20^\circ$. Then $\triangle BCG$ is similar to $\triangle DCE$ and $BC=2(EC)$. Also $\triangle FGC$ is congruent to $\triangle ABC$. Therefore, $[\triangle ACF]  = ([\triangle ABC] + [\triangle GCF]) + [\triangle BCG]$. Plugging in the values that we know and then dividing by 2 results in an answer of $\boxed{B) \frac{\sqrt{3}}{8}.}$

 This solution is from the solution manual but was typed here by alpha\_2.

 


\end{enumerate}